\begin{titlepage}
  \thispagestyle{empty}
  \vspace*{3cm}
  {\sffamily\color{inkgray}\large LiteBee Wing \quad·\quad LiteBee Bot (C10) \quad·\quad модуль позиционирования\par}
  \vspace{1.2cm}
  {\sffamily\bfseries\color{accent}\fontsize{34}{40}\selectfont Как устроены\\дроны LiteBee\par}
  \vspace{0.8cm}
  {\sffamily\color{ink}\Large От метки на ковре до мотора:\\разбор прошивки и протоколов\par}
  \vspace{1.4cm}
  \begin{tikzpicture}
    \draw[accent, line width=2pt] (0,0) -- (3.2,0);
  \end{tikzpicture}
  \vfill
  {\sffamily\color{inkgray}\small
  Книга собрана из разбора прошивок производителя (образ дрона Wing, ядро полёта C10, мобильные приложения), из файлов, снятых с самого дрона и с модуля позиционирования по USB, и из живых измерений на дроне 6 ноября 2026 года.\par
  \medskip
  Каждое утверждение помечено: \pr{} (видно в файле или на устройстве), \vyv{} (наше заключение) или \nei{}. Исходные тексты разбора лежат в папке \texttt{docs/firmware}, скрипты в \texttt{tools/firmware}.\par}
\end{titlepage}
